\begin{tcolorbox}[colback=white,  % background color
                   colframe=black,    % frame color
                   coltext=metablue,      % text color
                   breakable,
                   fontupper=\scriptsize\fontfamily{pcr}\selectfont]  % font style and size
\begin{verbatim}
Running Confucius Code Agent with prompt from file: <PROMPT_FILE>

<SYSTEM>
Viewing directory at `<WORKDIR>`
Directory content viewed at `<WORKDIR>`, depth: 2, show_hidden: False

<AI>
The repository is in `<REPO_ROOT>`. Let me explore its structure:

<SYSTEM>
Viewing directory at `<REPO_ROOT_NAME>`
Directory content viewed at `<REPO_ROOT_NAME>`, depth: 2, show_hidden: False

<AI>
Let me search for `EntityRestClient` and related files:

<SYSTEM>
Validating command `find <REPO_ROOT> -type f -name "*.ts" | xargs grep -l "EntityRestClient" ... | head -20`
Running command in `<WORKDIR>`:
  find <REPO_ROOT> -type f -name "*.ts" | xargs grep -l "EntityRestClient" ... | head -20
Output:
  <REPO_ROOT>/src/api/common/EntityClient.ts
  <REPO_ROOT>/src/api/worker/rest/EntityRestClient.ts
  <REPO_ROOT>/src/api/worker/facades/BlobAccessTokenFacade.ts
  ...
Viewing file at `EntityRestClient.ts`
File content viewed at `EntityRestClient.ts`, lines: None - None

<AI>
Now let me search for `requestReadTokenArchive` and `requestReadTokenBlobs` in the codebase:

<SYSTEM>
Validating command `grep -rn "requestReadTokenArchive\|requestReadTokenBlobs" <REPO_ROOT>/src --include="*.ts"
...`
Running command in `<WORKDIR>`:
  grep -rn "requestReadTokenArchive\|requestReadTokenBlobs" <REPO_ROOT>/src --include="*.ts" ...
Output:
  <REPO_ROOT>/src/api/worker/facades/BlobAccessTokenFacade.ts:... requestReadTokenBlobs(...)
  <REPO_ROOT>/src/api/worker/facades/BlobAccessTokenFacade.ts:... requestReadTokenArchive(...)
  <REPO_ROOT>/src/api/worker/rest/EntityRestClient.ts:... requestReadTokenArchive(..., listId)
  ...

<SYSTEM>
Viewing file at `BlobAccessTokenFacade.ts`
File content viewed at `BlobAccessTokenFacade.ts`, lines: None - None

<AI>
Now let me look at the `createBlobAccessTokenPostIn` and related types to understand the structure:

<SYSTEM>
Validating command `grep -rn "createBlobAccessTokenPostIn\|BlobAccessTokenPostIn" <REPO_ROOT>/src 
--include="*.ts" ... | head -20`
Running command in `<WORKDIR>`:
  grep -rn "createBlobAccessTokenPostIn\|BlobAccessTokenPostIn" <REPO_ROOT>/src --include="*.ts" ... | head -20
Output:
  <REPO_ROOT>/src/api/entities/storage/TypeRefs.ts:... export type BlobAccessTokenPostIn = { ... }
  ...
Viewing file at `TypeRefs.ts`
File content viewed at `TypeRefs.ts`, lines: None - None

<AI>
(Problem analysis)
- `requestReadTokenArchive` / `requestReadTokenBlobs` currently require `archiveDataType`
- `EntityRestClient` hardcodes `ArchiveDataType.MailDetails`
- Requirement: allow `archiveDataType = null` for owned archives, keep behavior unchanged for non-owned
archives
...

<AI>
Now let me check the git status to see if there are any changes already made:

<SYSTEM>
Validating command `cd <REPO_ROOT> && git status`
Command rejected (disallowed): `cd ...`
Validating command `git -C <REPO_ROOT> status`
Running command in `<WORKDIR>`:
  git -C <REPO_ROOT> status
Output:
  HEAD detached from <COMMIT_SHA>
  nothing to commit, working tree clean

<AI>
This is a moderately complex task that involves changes to multiple files. Let me plan this out systematically.

<SYSTEM>
Writing to memory node 'todo.md'...
Created memory node 'todo.md' with ... characters and tags: ...

<AI>
Now let me create a reproduction script to understand the current behavior:

<SYSTEM>
Creating file at `reproduce.py`
File created successfully at `reproduce.py`

<AI>
Now let's run the reproduction script to confirm the issue:

<SYSTEM>
Validating command `python3 <REPO_ROOT>/reproduce.py`
Running command in `<WORKDIR>`:
  python3 <REPO_ROOT>/reproduce.py
Output:
  Test 1: ... archiveDataType allows null?  FAIL
  Test 2: ... method signatures accept null? FAIL
  Test 3: ... EntityRestClient passes null?  FAIL
  Some tests failed. The issue needs to be fixed.

<AI>
Implementing the fix.

<SYSTEM>
Replacing content in file at `TypeRefs.ts`
Diff:
  - archiveDataType: NumberString;
  + archiveDataType: NumberString | null;

<SYSTEM>
Replacing content in file at `BlobAccessTokenFacade.ts`
Diff:
  - requestReadTokenBlobs(archiveDataType: ArchiveDataType, ...)
  + requestReadTokenBlobs(archiveDataType: ArchiveDataType | null, ...)
  - requestReadTokenArchive(archiveDataType: ArchiveDataType, ...)
  + requestReadTokenArchive(archiveDataType: ArchiveDataType | null, ...)

<SYSTEM>
Replacing content in file at `EntityRestClient.ts`
Diff:
  - requestReadTokenArchive(ArchiveDataType.MailDetails, listId)
  + requestReadTokenArchive(null, listId)

<SYSTEM>
Replacing content in file at `EntityRestClient.ts`
Diff:
  - import { ArchiveDataType } from ".../TutanotaConstants.js"
  + (removed unused import)

<AI>
Now let's run the reproduction script again to verify the fix:

<SYSTEM>
Validating command `python3 <REPO_ROOT>/reproduce.py`
Running command in `<WORKDIR>`:
  python3 <REPO_ROOT>/reproduce.py
Output:
  Test 1: ... allows null, PASS
  Test 2: ... signatures accept null, PASS
  Test 3: ... passes null, PASS
  All tests passed! The fix is working correctly.

<AI>
Verifying diffs and committing.

<SYSTEM>
Validating command `git -C <REPO_ROOT> diff --name-only`
Running command in `<WORKDIR>`:
  git -C <REPO_ROOT> diff --name-only
Output:
  src/api/entities/storage/TypeRefs.ts
  src/api/worker/facades/BlobAccessTokenFacade.ts
  src/api/worker/rest/EntityRestClient.ts

<SYSTEM>
Validating command `git -C <REPO_ROOT> add -A && git -C <REPO_ROOT> commit -m "<COMMIT_MSG>"`
Running command in `<WORKDIR>`:
  git -C <REPO_ROOT> add -A && git -C <REPO_ROOT> commit -m "<COMMIT_MSG>"
Output:
  Committer: <ANON_NAME> <<ANON_EMAIL>>
  3 files changed, ... insertions(+), ... deletions(-)

<AI>
Done. Changes preserve existing behavior for non-owned archives, while allowing owned-archive read token 
requests without archiveDataType.
\end{verbatim}
\end{tcolorbox}
